\documentclass[12pt,a4paper]{article}
\usepackage[none]{hyphenat}
\title{Maths is Fun: Soapy Solutions}
\author{Disha Kuzhively}
\begin{document}
\maketitle
\begin{abstract}
    Finding the least, the fastest, or the most efficient way to do something is a central idea that appears everywhere,  from everyday decisions to advanced mathematics. Nature does this beautifully when a soap film stretches across a wire frame, it naturally settles into a surface of least area. But this raises a question. Given any boundary curve no matter how irregular, does such a ``smallest possible surface" always exist?

    In this hands-on workshop, we will explore this question using soap films. We will dip, twist, and experiment with different wire frames, and watch how surfaces form in response. There are no formulas required. Just curiosity, careful observation, and a willingness to get a little soapy.
\end{abstract}
\end{document}